\documentclass[UTF8,fontset=fandol,zihao=-4,a4paper]{ctexart}

\usepackage[top=2.25cm,bottom=2.20cm,left=2.35cm,right=2.35cm]{geometry}
\usepackage{fontspec}
\setmonofont{Menlo}[Scale=0.92]
\usepackage{amsmath,amssymb}
\usepackage{booktabs,array,tabularx}
\usepackage{graphicx,caption}
\usepackage{xcolor}
\usepackage{fancyhdr}
\usepackage{fvextra}
\usepackage{needspace}
\usepackage{microtype}

\definecolor{CodeText}{HTML}{1F2933}
\definecolor{CodeRule}{HTML}{A8B2BC}
\definecolor{FileBar}{HTML}{E9EEF2}

\pagestyle{fancy}
\fancyhf{}
\fancyfoot[C]{\thepage}
\renewcommand{\headrulewidth}{0pt}

\captionsetup{font=small,labelsep=quad}

% ---------- 代码排版（较范本放大字号，保证清晰可读） ----------
\newcommand{\filedivider}[1]{%
  \Needspace{3\baselineskip}%
  \par\vspace{0.6\baselineskip}%
  \noindent\colorbox{FileBar}{%
    \parbox{\dimexpr\linewidth-2\fboxsep\relax}{\small\ttfamily\detokenize{#1}}%
  }%
  \par\vspace{0.3\baselineskip}%
}

\newcommand{\codeinput}[1]{%
  \VerbatimInput[
    fontsize=\small,
    formatcom=\color{CodeText},
    numbers=left,
    numbersep=5pt,
    xleftmargin=2.4em,
    frame=single,
    framesep=3pt,
    rulecolor=\color{CodeRule},
    tabsize=4,
    breaklines=true,
    breakanywhere=false,
    breaksymbolleft={},
    breaksymbolright={},
    showspaces=false,
    showtabs=false
  ]{#1}%
}

% ---------- 参考文献条目 ----------
\newcommand{\refitem}[2]{\par\noindent\hangindent=2.2em\hangafter=1 [#1]\hspace{0.6em}#2\par\vspace{2pt}}

\begin{document}
\setcounter{page}{38}

\section*{参考文献}
\addcontentsline{toc}{section}{参考文献}

\refitem{1}{Lo Y M D, Corbetta N, Chamberlain P F, et al. Presence of fetal DNA in maternal plasma and serum[J]. The Lancet, 1997, 350(9076): 485--487.}
\refitem{2}{Chiu R W K, Chan K C A, Gao Y, et al. Noninvasive prenatal diagnosis of fetal chromosomal aneuploidy by massively parallel genomic sequencing of DNA in maternal plasma[J]. Proceedings of the National Academy of Sciences, 2008, 105(51): 20458--20463.}
\refitem{3}{Wang E, Batey A, Struble C, et al. Gestational age and maternal weight effects on fetal cell-free DNA in maternal plasma[J]. Prenatal Diagnosis, 2013, 33(7): 662--666.}
\refitem{4}{Ashoor G, Syngelaki A, Poon L C Y, et al. Fetal fraction in maternal plasma cell-free DNA at 11--13 weeks' gestation: relation to maternal and fetal characteristics[J]. Ultrasound in Obstetrics \& Gynecology, 2013, 41(1): 26--32.}
\refitem{5}{Bianchi D W, Parker R L, Wentworth J, et al. DNA sequencing versus standard prenatal aneuploidy screening[J]. New England Journal of Medicine, 2014, 370(9): 799--808.}
\refitem{6}{Norton M E, Jacobsson B, Swamy G K, et al. Cell-free DNA analysis for noninvasive examination of trisomy[J]. New England Journal of Medicine, 2015, 372(17): 1589--1597.}
\refitem{7}{Klein J P, Moeschberger M L. Survival Analysis: Techniques for Censored and Truncated Data[M]. 2nd ed. New York: Springer, 2003.}
\refitem{8}{王淑娟, 高志英, 吕艳萍, 等. 孕妇外周血中游离胎儿DNA检测在诊断胎儿染色体异常中的应用价值[J]. 中华妇产科杂志, 2012, 47(11): 808--812.}
\refitem{9}{赵军红, 代鹏, 朱若男, 等. 2398例孕妇外周血游离胎儿DNA产前筛查阳性结果的验证与分析[J]. 中华妇产科杂志, 2020, 55(10): 679--684.}
\refitem{10}{国家卫生计生委办公厅. 孕妇外周血胎儿游离DNA产前筛查与诊断技术规范[S]. 国卫办妇幼发〔2016〕45号, 2016.}

\clearpage
\section*{附录}
\addcontentsline{toc}{section}{附录}

\subsection*{附录 A\hspace{1em}补充数据表}
\addcontentsline{toc}{subsection}{附录 A 补充数据表}
\renewcommand{\thetable}{A.\arabic{table}}

本部分汇总正文中受篇幅限制未完整列出的中间结果。表~\ref{tab:a1}~为问题一候选混合效应模型的信息准则比较；表~\ref{tab:a2}~给出问题二不同分组数下的动态规划决策损失；表~\ref{tab:a3}~为问题二达标定义与技术误差设定的单因素敏感性；表~\ref{tab:a4}~展示问题三检测失败情景下的分组推荐时点；表~\ref{tab:a5}~与表~\ref{tab:a6}~分别给出问题四的阳性标签审计与候选特征集的折外性能。

\begin{table}[!htbp]
  \centering
  \caption{问题一候选混合效应模型的对数似然与信息准则（按 BIC 升序）}
  \label{tab:a1}
  \small
  \begin{tabular}{@{}lccc@{}}
    \toprule
    候选模型结构 & 对数似然 & AIC & BIC \\
    \midrule
    变点 18 周 & -393.73 & 801.45 & 835.96 \\
    变点 17 周 & -396.11 & 806.23 & 840.74 \\
    变点 16 周 & -398.87 & 811.74 & 846.24 \\
    线性（无变点） & -403.28 & 818.56 & 848.14 \\
    孕周 $\times$ BMI 交互 & -401.16 & 816.33 & 850.83 \\
    变点 15 周 & -401.23 & 816.46 & 850.97 \\
    变点 14 周 & -402.61 & 819.23 & 853.74 \\
    变点 13 周 & -403.16 & 820.32 & 854.82 \\
    加入年龄、身高 & -400.19 & 816.38 & 855.82 \\
    \bottomrule
  \end{tabular}
\end{table}


\begin{table}[!htbp]
  \centering
  \caption{问题二分组数 $K$ 的动态规划总损失与 BMI 切点（普通判定口径）}
  \label{tab:a2}
  \small
  \begin{tabular}{@{}ccccc@{}}
    \toprule
    可靠度 $\rho$ & 分组数 $K$ & 总损失 & 损失下降占比 & BMI 切点 \\
    \midrule
    0.80 & 1 & 122.19 & 0.0\% & --- \\
    0.80 & 2 & 74.67 & 49.7\% & 31.77 \\
    0.80 & 3 & 46.95 & 78.7\% & 31.03, 35.03 \\
    0.80 & 4 & 36.66 & 89.5\% & 30.91, 33.62, 36.10 \\
    0.80 & 5 & 30.97 & 95.4\% & 29.53, 31.03, 33.62, 36.10 \\
    0.80 & 6 & 26.62 & 100.0\% & 29.53, 30.86, 32.70, 34.26, 36.10 \\
    0.90 & 1 & 112.03 & 0.0\% & --- \\
    0.90 & 2 & 59.85 & 58.5\% & 33.81 \\
    0.90 & 3 & 38.88 & 82.0\% & 31.03, 34.82 \\
    0.90 & 4 & 29.27 & 92.7\% & 30.76, 33.62, 36.10 \\
    0.90 & 5 & 25.73 & 96.7\% & 29.93, 31.43, 33.62, 36.10 \\
    0.90 & 6 & 22.77 & 100.0\% & 29.27, 30.64, 32.19, 33.62, 36.10 \\
    \bottomrule
  \end{tabular}
\end{table}


\begin{table}[!htbp]
  \centering
  \caption{问题二达标阈值、可信水平与误差倍数的单因素敏感性（$\rho=0.80$）}
  \label{tab:a3}
  \small
  \begin{tabular}{@{}ccccccc@{}}
    \toprule
    判定方式 & 达标阈值 & 可信水平 $\eta$ & 误差倍数 & 最优分布 & BMI 系数 & 单位 BMI 时间比 \\
    \midrule
    硬判定 & 0.035 & 0.025 & 1.00 & 对数 logistic & 0.0445 & 1.0455 \\
    可信判定 & 0.035 & 0.025 & 1.00 & 对数 logistic & 0.1403 & 1.1506 \\
    硬判定 & 0.040 & 0.025 & 1.00 & 对数正态 & 0.0392 & 1.0400 \\
    可信判定 & 0.040 & 0.025 & 1.00 & 对数 logistic & 0.0830 & 1.0865 \\
    硬判定 & 0.045 & 0.025 & 1.00 & 对数正态 & 0.0365 & 1.0372 \\
    可信判定 & 0.045 & 0.025 & 1.00 & Weibull & 0.0520 & 1.0533 \\
    可信判定 & 0.040 & 0.010 & 1.00 & 对数 logistic & 0.1058 & 1.1116 \\
    可信判定 & 0.040 & 0.050 & 1.00 & 对数 logistic & 0.0746 & 1.0775 \\
    可信判定 & 0.040 & 0.100 & 1.00 & Weibull & 0.0484 & 1.0496 \\
    可信判定 & 0.040 & 0.025 & 0.75 & Weibull & 0.0560 & 1.0576 \\
    可信判定 & 0.040 & 0.025 & 1.25 & 对数 logistic & 0.1041 & 1.1098 \\
    可信判定 & 0.040 & 0.025 & 1.50 & 对数 logistic & 0.1047 & 1.1104 \\
    \bottomrule
  \end{tabular}
\end{table}


\begin{table}[!htbp]
  \centering
  \caption{问题三检测失败概率与复测延迟情景下的分组推荐时点}
  \label{tab:a4}
  \small
  \begin{tabular}{@{}ccccccc@{}}
    \toprule
    可靠度 $\rho$ & 失败概率 & 复测延迟（周） & 第 1 组 & 第 2 组 & 第 3 组 & 第 4 组 \\
    \midrule
    0.80 & 0\% & 1 & 11周+4天 & 12周+6天 & 14周+0天 & 16周+2天 \\
    0.80 & 0\% & 2 & 11周+4天 & 12周+6天 & 14周+0天 & 16周+2天 \\
    0.80 & 5\% & 1 & 11周+4天 & 12周+6天 & 14周+0天 & 16周+3天 \\
    0.80 & 5\% & 2 & 11周+5天 & 13周+0天 & 14周+1天 & 16周+3天 \\
    0.80 & 10\% & 1 & 11周+5天 & 13周+0天 & 14周+1天 & 16周+3天 \\
    0.80 & 10\% & 2 & 11周+6天 & 13周+0天 & 14周+2天 & 16周+4天 \\
    0.80 & 15\% & 1 & 11周+5天 & 13周+0天 & 14周+1天 & 16周+3天 \\
    0.80 & 15\% & 2 & 11周+6天 & 13周+1天 & 14周+2天 & 16周+5天 \\
    0.90 & 0\% & 1 & 14周+5天 & 16周+2天 & 17周+6天 & 22周+1天 \\
    0.90 & 0\% & 2 & 14周+5天 & 16周+2天 & 17周+6天 & 22周+1天 \\
    0.90 & 5\% & 1 & 14周+5天 & 16周+3天 & 18周+0天 & 22周+2天 \\
    0.90 & 5\% & 2 & 14周+6天 & 16周+3天 & 18周+0天 & 22周+2天 \\
    0.90 & 10\% & 1 & 14周+6天 & 16周+3天 & 18周+0天 & 22周+2天 \\
    0.90 & 10\% & 2 & 15周+0天 & 16周+4天 & 18周+1天 & 22周+3天 \\
    0.90 & 15\% & 1 & 14周+6天 & 16周+4天 & 18周+0天 & 22周+2天 \\
    0.90 & 15\% & 2 & 15周+0天 & 16周+5天 & 18周+2天 & 22周+4天 \\
    \bottomrule
  \end{tabular}
\end{table}


\begin{table}[!htbp]
  \centering
  \caption{问题四女胎阳性标签规模与 $|Z|\geq 3$ 规则覆盖审计}
  \label{tab:a5}
  \small
  \begin{tabular}{@{}ccccc@{}}
    \toprule
    异常类型 & 阳性记录数 & 阳性孕妇数 & $|Z|\geq 3$ 命中比例 & 阳性 $Z$ 值范围 \\
    \midrule
    T13 & 23 & 18 & 4.3\% & [-2.81, 3.81] \\
    T18 & 46 & 30 & 0.0\% & [-0.90, 2.97] \\
    T21 & 13 & 12 & 0.0\% & [-2.39, 1.33] \\
    \bottomrule
  \end{tabular}
\end{table}


\begin{table}[!htbp]
  \centering
  \caption{问题四候选特征集折外性能（多种子加权均值，弹性网络逻辑回归）}
  \label{tab:a6}
  \small
  \begin{tabular}{@{}llccccc@{}}
    \toprule
    任务 & 特征集 & 特征数 & ROC-AUC & PR-AUC & F1 & F1 的 95\% 区间 \\
    \midrule
    ANY & 仅 Z 值 & 9 & 0.485 & 0.123 & 0.193 & [0.169, 0.207] \\
    ANY & 核心测序特征 & 29 & 0.784 & 0.435 & 0.426 & [0.377, 0.459] \\
    ANY & 核心+手工特征 & 64 & 0.743 & 0.387 & 0.412 & [0.337, 0.477] \\
    ANY & 加入纵向特征 & 157 & 0.703 & 0.308 & 0.325 & [0.307, 0.355] \\
    ANY & 去除孕妇特征 & 149 & 0.686 & 0.320 & 0.327 & [0.275, 0.390] \\
    ANY & 去除质控特征 & 87 & 0.675 & 0.297 & 0.292 & [0.252, 0.345] \\
    T13 & 仅 Z 值 & 9 & 0.367 & 0.034 & 0.065 & [0.054, 0.075] \\
    T13 & 核心测序特征 & 29 & 0.823 & 0.230 & 0.312 & [0.237, 0.374] \\
    T13 & 核心+手工特征 & 64 & 0.807 & 0.184 & 0.234 & [0.134, 0.272] \\
    T13 & 加入纵向特征 & 157 & 0.828 & 0.210 & 0.270 & [0.187, 0.349] \\
    T13 & 去除孕妇特征 & 149 & 0.846 & 0.216 & 0.272 & [0.237, 0.343] \\
    T13 & 去除质控特征 & 87 & 0.643 & 0.075 & 0.102 & [0.043, 0.162] \\
    T18 & 仅 Z 值 & 9 & 0.571 & 0.094 & 0.161 & [0.130, 0.185] \\
    T18 & 核心测序特征 & 29 & 0.850 & 0.517 & 0.461 & [0.383, 0.537] \\
    T18 & 核心+手工特征 & 64 & 0.830 & 0.475 & 0.471 & [0.422, 0.514] \\
    T18 & 加入纵向特征 & 157 & 0.789 & 0.379 & 0.386 & [0.299, 0.469] \\
    T18 & 去除孕妇特征 & 149 & 0.789 & 0.401 & 0.426 & [0.355, 0.516] \\
    T18 & 去除质控特征 & 87 & 0.723 & 0.285 & 0.341 & [0.268, 0.408] \\
    T21 & 仅 Z 值 & 9 & 0.334 & 0.018 & 0.042 & [0.033, 0.047] \\
    T21 & 核心测序特征 & 29 & 0.497 & 0.179 & 0.059 & [0.044, 0.091] \\
    T21 & 核心+手工特征 & 64 & 0.453 & 0.157 & 0.042 & [0.018, 0.054] \\
    T21 & 加入纵向特征 & 157 & 0.416 & 0.041 & 0.025 & [0.000, 0.050] \\
    T21 & 去除孕妇特征 & 149 & 0.387 & 0.035 & 0.020 & [0.000, 0.045] \\
    T21 & 去除质控特征 & 87 & 0.376 & 0.020 & 0.040 & [0.029, 0.047] \\
    \bottomrule
  \end{tabular}
\end{table}


\clearpage
\subsection*{附录 B\hspace{1em}核心源代码}
\addcontentsline{toc}{subsection}{附录 B 核心源代码}

本部分给出模型的核心求解程序。代码已脱敏处理，数据文件以相对路径``附件.xlsx''读入，不含任何本地绝对路径；关键步骤处添加了必要注释。运行环境为 Python 3.11 及以上版本，依赖 numpy、pandas、scipy、statsmodels 等开源库。固定随机种子 \texttt{SEED = 20250904}，全部结果可复现。

\filedivider{nipt\_core.py（核心函数库：数据准备、技术误差估计、区间删失 AFT、动态规划分组）}
\codeinput{code/nipt_core_annotated.py}

\filedivider{solve\_baseline.py（主流程：问题一至问题四的基线求解串联）}
\codeinput{code/solve_baseline_annotated.py}

\end{document}
